\subsubsection{Replacement Sort / Heap Sort}
It has the advantage that it can produce longer runs with the same amount of memory, on average $2B$ (instead of $B$) with $B$ frames of memory.
It is used as the sorting algorithm in pass $0$ above.

It works as follows:
\begin{algorithm}
    \caption{Replacement Sort}
    \begin{algorithmic}[1]
        \Procedure{ReplacementSort}{$R$}
            \State Read $B - 1$ pages of relation $R$, store in a heap (priority queue)
            \For{every page}
                \If{Top of heap is smaller than the end of the sorted run}
                    \State Continue with next iteration
                \ElsIf{No element is left in memory that is larger than last element of current sorted run}
                    \State create new run
                \Else
                    \State load page, remove smaller records and place them on the current sorted run, until a frame is freed.
                \EndIf
            \EndFor
        \EndProcedure
    \end{algorithmic}
\end{algorithm}

Informally, it is searching for an always larger value once it has committed to a certain value.
It picks the first value to be the lowest value in the current frames.

A new run is started if all values in the loaded frames are deferred or used already.

This means that after the sort run is complete, a merge is still needed.

\paragraph{Performance}
\bi{Best case}: $N$ elements in a run (thus a single run, data is already sorted)

\bi{Worst case}: $B - 1$ elements in a run (reversed data)

\bi{Average case}: $2B$ elements in a run.

In any case, the cost of sorting is \cost{$2N(1 + \lceil \log_{B - 1} \lceil N / E \rceil \rceil)$}, with $E$ the number of elements ($2B$ in the average case)

Thus, Quicksort is faster, but longer runs often mean fewer passes, saving time.
